\documentclass[12pt,reqno]{amsart}

\usepackage{graphicx}

\usepackage{amssymb}
\usepackage{amsthm}
\theoremstyle{plain}

\newtheorem*{theorem*}{Theorem}
%% this allows for theorems which are not automatically numbered

\newtheorem{definition}{Definition}
\newtheorem{theorem}{Theorem}
\newtheorem{lemma}{Lemma}
\newtheorem{example}{Example}
\newtheorem{proposition}{Proposition}
\newtheorem{corollary}{Corollary}
\usepackage{lineno}

%% The above lines are for formatting.  In general, you will not want to change these.


\title{Exact Ornstein-Uhlenbeck Numerical Update Rule}

\author{Luca Roncoroni Seda}

\begin{document}

\begin{abstract}
    In this supplement, we derive the exact numerical update rule used in simulating the Ornstein-Uhlenbeck magnetic field model, justifying its usage as opposed to the Euler-Maruyama scheme.
\end{abstract}

\maketitle

Let $X_t$ be a random variable satisfying the Ornstein-Uhlenbeck SDE
\begin{equation*}
    dX_t = -\lambda X_t dt + \sigma dW_t, \quad \lambda = \frac{1}{\tau_c} > 0.
\end{equation*}
We wish to derive an exact discrete-time update for $X_{n+1} = X(t_n+\Delta t)$ given $X_n = X(t_n)$.

\section{Existence and form of the solution}
The SDE is linear with globally Lipschitz drift $b(x) = -\lambda x$ and constant diffusion coefficient $\sigma(x) = \sigma$. Therefore, by Itô's existence and uniqueness theorem, there exists a unique strong solution to the SDE.

\begin{theorem}{(Itô's product rule)}
    If $X_t$ and $Y_t$ are semimartingales, then
    \begin{equation*}
        d(X_t Y_t) = X_{t} dY_t + Y_{t} dX_t + d[X,Y]_t,
    \end{equation*}
    where $[X,Y]_t$ is the quadratic covariation of $X_t$ and $Y_t$.
    \label{thm:product_rule}
\end{theorem}

\begin{proposition}{}
    The discrete-time solution to the Ornstein-Uhlenbeck SDE is given by
    \begin{equation*}
        X_{n+1} = e^{-\lambda \Delta t} X_n + \sigma \int_{t_n}^{t_{n+1}} e^{-\lambda (t_{n+1}-s)} dW_s.
    \end{equation*}
    \label{prop:discrete_solution}
\end{proposition}
\begin{proof}
    Using the integrating factor $e^{\lambda t}$, a finite-variation process satisfying $d[e^{\lambda t}, X_t] = 0$, by Itô's product rule we have
    \begin{align*}
        d(e^{\lambda t} X_t) &= e^{\lambda t}dX_t + \lambda e^{\lambda t} X_t dt\\
        &= e^{\lambda t}(-\lambda X_t dt + \sigma dW_t) + \lambda e^{\lambda t} X_t dt\\
        &= e^{\lambda t} \sigma dW_t.
    \end{align*}
    Therefore, integrating from $t_n$ to $t_{n+1} = t_n + \Delta t$ gives
    \begin{align*}
        e^{\lambda t_{n+1}}X_{n+1} - e^{\lambda t_n}X_n &= \sigma \int_{t_n}^{t_{n+1}} e^{\lambda s} dW_s\\
        X_{n+1} &= e^{-\lambda \Delta t} X_n + \sigma \int_{t_n}^{t_{n+1}} e^{-\lambda (t_{n+1}-s)} dW_s.
    \end{align*}
\end{proof}


\section{Distribution of the stochastic integral}
\begin{theorem}{(Itô isometry)}
    Let $W_t$ be a Brownian motion and let $f\in L^2([a,b])$ be deterministic. Then
    \begin{equation*}
        \mathbb{E}\left[\left(\int_a^b f(s) dW_s\right)^2\right] = \int_a^b f(s)^2 ds.
    \end{equation*}
    In particular,
    \begin{equation*}
        \int_a^b f(s) dW_s \sim \mathcal{N}\left(0, \int_a^b f(s)^2 ds\right).
    \end{equation*}
    \label{thm:ito_isometry}
\end{theorem}

\begin{corollary}{}
    Let $X_t$ satisfy the Ornstein-Uhlenbeck SDE. Then, for any timestep $\Delta t > 0$,
    \begin{equation*}
        X_{n+1}\,\vert\,X_n \sim \mathcal{N}\left( e^{-\lambda \Delta t} X_n, \sigma^2 \int_{t_n}^{t_{n+1}} e^{-2\lambda (t_{n+1}-s)} ds \right).
    \end{equation*}
    Equivalently,
    \begin{equation}
        X_{n+1} = e^{-\lambda \Delta t} X_n + \sigma \sqrt{\frac{1-e^{-2\lambda \Delta t}}{2\lambda}}\,\xi_n,\quad \xi_n \overset{\mathrm{iid}}{\sim} \mathcal{N}(0,1)
    \end{equation}
    \label{cor:discrete_distribution}
\end{corollary}

\begin{proof}
    By Proposition~\ref{prop:discrete_solution},
    \begin{equation*}
        X_{n+1} = e^{-\lambda \Delta t} X_n + \sigma \int_{t_n}^{t_{n+1}} e^{-\lambda (t_{n+1}-s)} dW_s.
    \end{equation*}
    The stochastic integral is Gaussian with mean zero. By Itô isometry,
    \begin{align*}
        \mathrm{Var}\,\left[\int_{t_n}^{t_{n+1}} e^{-\lambda (t_{n+1}-s)} dW_s\right] &= \int_{t_n}^{t_{n+1}} e^{-2\lambda (t_{n+1}-s)} ds\\
        &= \frac{1-e^{-2\lambda\Delta t}}{2\lambda}.
    \end{align*}
    The result follows.
\end{proof}

Letting $\lambda = 1/\tau_c$, we obtain the form encountered in the main text:
\begin{equation*}
    X_{n+1} = e^{-\Delta t/\tau_c} X_n + \sigma \sqrt{\frac{\tau_c}{2}\left(1-e^{-2\Delta t/\tau_c}\right)}\,\xi_n.
\end{equation*}
Furthermore, note that the $\xi_n$ are independent because the stochastic integrals are taken over disjoint Brownian intervals.

\section{Preference over Euler-Maruyama}
Since the Ornstein–Uhlenbeck process possesses a closed-form solution, its transition law may be sampled exactly. This preserves the prescribed variance and exponential correlation function for arbitrary timestep and avoids the timestep-dependent bias introduced by Euler–Maruyama discretization. The exact update is therefore both more accurate and computationally more efficient for the correlation times considered in this work.


\end{document}